\documentclass[11pt,a4paper,UTF8]{ctexart}
\usepackage{geometry}
\geometry{left=18mm,right=18mm,top=24mm,bottom=24mm,headheight=22pt,footskip=12mm}
\usepackage{amsmath,amssymb,mathtools,bm,esint,extarrows}
\usepackage{xcolor}
\usepackage{tcolorbox}
\tcbuselibrary{skins,breakable}
\usepackage{fancyhdr}
\usepackage{titlesec}
\usepackage{hyperref}
\usepackage{tikz}
\usepackage{booktabs}
\usepackage{tabularx}

% --- Color Palette (Purple & DeepPink Academic Theme) ---
\definecolor{DeepPurple}{HTML}{4C1D95}
\definecolor{TitlePurple}{HTML}{581C87}
\definecolor{SubPurple}{HTML}{6B21A8}
\definecolor{DeepPink}{HTML}{FF1493}
\definecolor{LightPinkBg}{HTML}{FFF5F8}
\definecolor{LightPurpleBg}{HTML}{F8F5FF}
\definecolor{GoldAccent}{HTML}{D97706}
\definecolor{DarkText}{HTML}{1F2937}

\hypersetup{
    colorlinks=true,
    linkcolor=TitlePurple,
    citecolor=DeepPink,
    urlcolor=SubPurple,
    pdfborder={0 0 0}
}

% --- Typography & Spacing ---
\linespread{1.3}
\setlength{\parskip}{0.35em plus 0.1em minus 0.1em}
\setlength{\parindent}{0pt}

% --- Section Titles Styling ---
\titleformat{\section}
  {\Large\bfseries\color{TitlePurple}}{\thesection}{1em}{}
  [\vspace{1mm}{\color{TitlePurple!30}\hrule height 1pt}]
\titleformat{\subsection}
  {\large\bfseries\color{SubPurple}}{\thesubsection}{0.8em}{}
\titleformat{\subsubsection}
  {\normalsize\bfseries\color{DeepPurple}}{\thesubsubsection}{0.6em}{}

% --- tcolorbox Environments ---
% 1. Theorem / Formula / Method / Analysis Box (Deep Purple)
\newtcolorbox{thmbox}[1][]{
    enhanced,
    breakable,
    colback=LightPurpleBg,
    colframe=TitlePurple,
    arc=3pt,
    boxrule=0pt,
    leftrule=3.5pt,
    left=10pt,
    right=10pt,
    top=8pt,
    bottom=8pt,
    title=#1,
    coltitle=white,
    colbacktitle=TitlePurple,
    fonttitle=\bfseries\small,
    attach boxed title to top left={yshift=-2mm, xshift=4mm},
    boxed title style={boxrule=0pt, arc=2pt}
}

% 2. Solution / Formula / Derivation Box (DeepPink)
\newtcolorbox{solbox}[1][]{
    enhanced,
    breakable,
    colback=LightPinkBg,
    colframe=DeepPink,
    arc=3pt,
    boxrule=0pt,
    leftrule=3.5pt,
    left=10pt,
    right=10pt,
    top=8pt,
    bottom=8pt,
    title=#1,
    coltitle=white,
    colbacktitle=DeepPink,
    fonttitle=\bfseries\small,
    attach boxed title to top left={yshift=-2mm, xshift=4mm},
    boxed title style={boxrule=0pt, arc=2pt}
}

% 3. Learning Goals / Summary Box
\newtcolorbox{targetbox}[1][]{
    enhanced,
    breakable,
    colback=LightPurpleBg,
    colframe=DeepPurple,
    arc=3pt,
    boxrule=1pt,
    left=10pt,
    right=10pt,
    top=8pt,
    bottom=8pt,
    title=#1,
    coltitle=white,
    colbacktitle=DeepPurple,
    fonttitle=\bfseries\small,
    attach boxed title to top left={yshift=-2mm, xshift=4mm},
    boxed title style={boxrule=0pt, arc=2pt}
}

\begin{document}

% ------------------- 讲义封面 (Cover Page) -------------------
\begin{titlepage}
\thispagestyle{empty}
\begin{tikzpicture}[remember picture,overlay]
    \fill[DeepPurple] (current page.north west) rectangle ([yshift=-7cm]current page.north east);
    \draw[DeepPink, line width=3.5pt] ([yshift=-7cm]current page.north west) -- ([yshift=-7cm]current page.north east);
    \draw[GoldAccent, line width=1pt] ([yshift=-7.12cm]current page.north west) -- ([yshift=-7.12cm]current page.north east);
    
    \node[anchor=north east, opacity=0.12, text=white] at ([xshift=-1.5cm, yshift=-0.5cm]current page.north east) {
        \fontsize{70}{70}\selectfont\bfseries 2027
    };
\end{tikzpicture}

\vspace*{0.8cm}
\begin{center}
    {\color{white}\fontsize{26}{32}\selectfont\bfseries 2027 考研数学(二) 核心讲义}\\[0.6cm]
    {\color{white}\fontsize{13}{18}\selectfont\bfseries 全国硕士研究生招生考试 · 统考数学(二) 核心精讲全书}\\[1.5cm]
\end{center}

\vspace{3.5cm}
\begin{center}
    {\color{GoldAccent}\large\bfseries $\blacktriangleright$ 基础强化 · 核心题解 · 考点深水区 $\blacktriangleleft$}\\[0.8cm]
    {\color{TitlePurple}\fontsize{24}{28}\selectfont\bfseries 一元积分学的应用（积分等式和积分不等式）· 核心例题精析}\\[0.6cm]
    {\color{DeepPink}\large\bfseries —— 核心精讲大例题与题解三步法 ——}\\[1.5cm]
\end{center}

\vfill

\begin{center}
\begin{tcolorbox}[
    colback=white,
    colframe=DeepPurple,
    width=0.88\textwidth,
    arc=4pt,
    boxrule=1.2pt,
    halign=center,
    drop shadow=gray!30
]
    \vspace{3mm}
    {\bfseries\large 编著团队：EscoffierZhou}\\[2.5mm]
    {\bfseries\color{TitlePurple} 目标院校：中国科学院大学 (UCAS) 计算机专硕 (22408)}\\[2.5mm]
    {\small\color{gray} 备考铁律：手算真题数字 · 错题白纸重做 · 408 客观题为王}\\[2.5mm]
    {\small 编制版本：2027 届首发版 · \today}
    \vspace{2mm}
\end{tcolorbox}
\end{center}
\vspace{0.8cm}
\end{titlepage}

% ------------------- 目录与页面主体配置 -------------------
\pagestyle{fancy}
\fancyhf{}
\fancyhead[L]{\small\color{gray}2027 考研数学(二) · 核心讲义系列}
\fancyhead[R]{\small\color{TitlePurple}\nouppercase{\leftmark}}
\fancyfoot[C]{\small\color{gray}— 第 \thepage\ 页 —}
\fancyfoot[R]{\small\color{gray}UCAS 22408}
\renewcommand{\headrulewidth}{0.5pt}

\pagenumbering{Roman}
\tableofcontents
\newpage
\pagenumbering{arabic}


---


\section{例题 11.1：推广积分第一中值定理与含参积分极限}


\subsection{1. 题目描述}
设函数 $f(x), g(x)$ 在闭区间 $[a, b]$ 上连续，且 $g(x)$ 在 $[a, b]$ 上不变号（即恒有 $g(x) \ge 0$ 或 $g(x) \le 0$）。
\textbf{(1)}  证明推广的积分第一中值定理：存在 $\xi \in [a, b]$，使得

\begin{equation*}
\int\_a^b f(x)g(x)\,\mathrm{d}x = f(\xi)\int\_a^b g(x)\,\mathrm{d}x
\end{equation*}

\textbf{(2)}  设 $f(x)$ 在 $[1, 2]$ 上连续，计算极限：

\begin{equation*}
\lim\_{n \to \infty} \int\_1^2 f(x)\mathrm{e}^{-x^n}\,\mathrm{d}x
\end{equation*}



\subsection{2. 解题思路与方法提炼}
\begin{solbox}[分步推导与答题规范]
\textbf{[1]}  \textbf{推广积分第一中值定理的双重视角}：
法一采用连续函数的\textbf{介值定理}。由闭区间上连续函数的极值定理，设 $m \le f(x) \le M$。两端同乘非负连续函数 $g(x)$ 并积分，由介值定理必存在介值点；法二将问题构造为变上限积分函数之商，直接应用\textbf{柯西中值定理}。
\textbf{[2]}  \textbf{积分极限的处理}：
对于含参变量积分极限，优先考查被积函数的变化趋势。在区间 $[1, 2]$ 上，当 $n \to \infty$ 时，$\mathrm{e}^{-x^n}$ 急剧衰减趋于 $0$。由闭区间连续函数的有界性，结合夹逼定理可直接证明极限为 $0$。
\end{solbox}


\subsection{3. 详细解答与规范步骤}


\subsubsection{(1) 推广积分第一中值定理的证明}
\begin{solbox}[分步推导与答题规范]
\textbf{[1]}  \textbf{确定函数界限}：
因为 $f(x)$ 在 $[a, b]$ 上连续，由闭区间连续函数的极值定理，存在最小值 $m$ 与最大值 $M$，使得对任意 $x \in [a, b]$ 均有：

\begin{equation*}
m \le f(x) \le M
\end{equation*}


\textbf{[2]}  \textbf{保持不等号同向积分}：
不妨设 $g(x) \ge 0$（若 $g(x) \le 0$，只需考查 $-g(x)$）。由于 $g(x) \ge 0$，用 $g(x)$ 乘以不等式各项：

\begin{equation*}
m g(x) \le f(x)g(x) \le M g(x)
\end{equation*}

在区间 $[a, b]$ 上逐项积分，由定积分保序性得：

\begin{equation*}
m \int\_a^b g(x)\,\mathrm{d}x \le \int\_a^b f(x)g(x)\,\mathrm{d}x \le M \int\_a^b g(x)\,\mathrm{d}x
\end{equation*}


\textbf{[3]}  \textbf{分类讨论并应用介值定理}：
情形一：若 $\int_a^b g(x)\,\mathrm{d}x = 0$。因 $g(x) \ge 0$ 连续，必有 $g(x) \equiv 0$。此时等式两边均为 $0$，对任意 $\xi \in [a, b]$ 等式均恒成立。\\ \\
情形二：若 $\int_a^b g(x)\,\mathrm{d}x > 0$，两边同除以此积分得：

\begin{equation*}
m \le \frac{\int\_a^b f(x)g(x)\,\mathrm{d}x}{\int\_a^b g(x)\,\mathrm{d}x} \le M
\end{equation*}

由闭区间连续函数的介值定理，必然存在 $\xi \in [a, b]$，使得：

\begin{equation*}
f(\xi) = \frac{\int\_a^b f(x)g(x)\,\mathrm{d}x}{\int\_a^b g(x)\,\mathrm{d}x}
\end{equation*}

即：

\begin{equation*}
\int\_a^b f(x)g(x)\,\mathrm{d}x = f(\xi)\int\_a^b g(x)\,\mathrm{d}x
\end{equation*}

\end{solbox}


\subsubsection{(2) 极限计算}
\begin{solbox}[分步推导与答题规范]
\textbf{[1]}  \textbf{放缩被积函数}：
因为 $f(x)$ 在 $[1, 2]$ 上连续，故存在正实数 $M > 0$，使得对任意 $x \in [1, 2]$ 均有 $|f(x)| \le M$。\\ \\
当 $x \in [1, 2]$ 时，$x^n \ge 1 + n(x-1)$ 或直接利用 $x \ge 1 \implies x^n \ge 1 + n(x-1)$，亦有 $x^n \ge 1$ 且 $x^n$ 随 $x$ 单调递增。\\ \\
更精细地，考虑积分：

\begin{equation*}
\left|\int\_1^2 f(x)\mathrm{e}^{-x^n}\,\mathrm{d}x\right| \le \int\_1^2 |f(x)|\mathrm{e}^{-x^n}\,\mathrm{d}x \le M \int\_1^2 \mathrm{e}^{-x^n}\,\mathrm{d}x
\end{equation*}


\textbf{[2]}  \textbf{变量代换或区间拆分}：
对任意小的固定常数 $\delta \in (0, 1)$，将区间拆分为 $[1, 1+\delta]$ 与 $[1+\delta, 2]$：

\begin{equation*}
\int\_1^2 \mathrm{e}^{-x^n}\,\mathrm{d}x = \int\_1^{1+\delta}\mathrm{e}^{-x^n}\,\mathrm{d}x + \int\_{1+\delta}^2 \mathrm{e}^{-x^n}\,\mathrm{d}x \le \delta \cdot 1 + (1-\delta)\mathrm{e}^{-(1+\delta)^n}
\end{equation*}

当 $n \to \infty$ 时，$(1+\delta)^n \to +\infty$，故 $\mathrm{e}^{-(1+\delta)^n} \to 0$。因此：

\begin{equation*}
\limsup\_{n\to\infty} \int\_1^2 \mathrm{e}^{-x^n}\,\mathrm{d}x \le \delta
\end{equation*}

令 $\delta \to 0^+$，即知 $\lim_{n\to\infty} \int_1^2 \mathrm{e}^{-x^n}\,\mathrm{d}x = 0$。\\ \\
由夹逼定理，原极限为：

\begin{equation*}
\lim\_{n \to \infty} \int\_1^2 f(x)\mathrm{e}^{-x^n}\,\mathrm{d}x = 0
\end{equation*}

\end{solbox}


\subsection{4. 考点延展与防坑警示}
\begin{solbox}[分步推导与答题规范]
\textbf{[1]}  推广第一积分中值定理中，\textit{}$g(x)$ 不变号是不可或缺的关键条件\textit{}。若 $g(x)$ 变号，正负积分抵消，中值定理结论完全可能失效。
\textbf{[2]}  若题目要求开区间 $\xi \in (a, b)$，需预先判定 $f(x)$ 是否为常数函数：若 $f(x)$ 恒为常数，开区间任一点均满足；若 $f(x)$ 非常数，则连续函数在闭区间的最小值与最大值不可能在全区间取满，积分商必然严格落在开区间 $(m, M)$ 内部，由介值定理保证 $\xi \in (a, b)$。
\end{solbox}

---


\section{例题 11.2：积分等式、积分中值定理与导数零点存在性}


\subsection{1. 题目描述}
设函数 $f(x)$ 在闭区间 $[0, 1]$ 上二阶可导，满足 $f(0) = f(1) = 0$，且满足积分等式条件：

\begin{equation*}
\int\_0^1 f(x)\,\mathrm{d}x > 0
\end{equation*}

证明：存在 $\xi \in (0, 1)$，使得 $f''(\xi) < 0$。


\subsection{2. 解题思路与方法提炼}
\begin{solbox}[分步推导与答题规范]
\textbf{[1]}  \textbf{提取函数在区间内的正值点}：
由积分等式（积分值大于零）结合积分第一中值定理，可知存在点 $c \in (0, 1)$ 使得 $f(c) > 0$。
\textbf{[2]}  \textbf{两次应用微分中值定理定位二阶导数}：
在区间 $[0, c]$ 与 $[c, 1]$ 上分别应用拉格朗日中值定理，获得一阶导数正负相异的两个点 $\eta_1, \eta_2$；再在 $[\eta_1, \eta_2]$ 上应用拉格朗日中值定理，即可确定 $f''(\xi) < 0$。
\end{solbox}


\subsection{3. 详细解答与规范步骤}
\begin{solbox}[分步推导与答题规范]
\textbf{[1]}  \textbf{获取区间内的一阶均值正点}：
因为 $f(x)$ 在 $[0, 1]$ 上连续，由积分第一中值定理，存在 $c \in (0, 1)$，使得：

\begin{equation*}
\int\_0^1 f(x)\,\mathrm{d}x = f(c)(1-0) = f(c)
\end{equation*}

已知 $\int_0^1 f(x)\,\mathrm{d}x > 0$，故有：

\begin{equation*}
f(c) > 0
\end{equation*}


\textbf{[2]}  \textbf{构造一阶导数符号差异}：
在区间 $[0, c]$ 上，函数 $f(x)$ 满足拉格朗日中值定理条件，存在 $\eta_1 \in (0, c)$，使得：

\begin{equation*}
f'(\eta\_1) = \frac{f(c) - f(0)}{c - 0} = \frac{f(c)}{c} > 0
\end{equation*}

在区间 $[c, 1]$ 上，同理存在 $\eta_2 \in (c, 1)$，使得：

\begin{equation*}
f'(\eta\_2) = \frac{f(1) - f(c)}{1 - c} = \frac{-f(c)}{1 - c} < 0
\end{equation*}

显然 $0 < \eta_1 < c < \eta_2 < 1$。

\textbf{[3]}  \textbf{应用拉格朗日中值定理判定二阶导数}：
函数 $f'(x)$ 在区间 $[\eta_1, \eta_2]$ 上满足拉格朗日中值定理条件，存在 $\xi \in (\eta_1, \eta_2) \subset (0, 1)$，使得：

\begin{equation*}
f''(\xi) = \frac{f'(\eta\_2) - f'(\eta\_1)}{\eta\_2 - \eta\_1}
\end{equation*}

由于 $\eta_2 > \eta_1$，分母 $\eta_2 - \eta_1 > 0$；而分子满足：

\begin{equation*}
f'(\eta\_2) < 0,\quad f'(\eta\_1) > 0 \implies f'(\eta\_2) - f'(\eta\_1) < 0
\end{equation*}

因此：

\begin{equation*}
f''(\xi) < 0
\end{equation*}

命题得证。
\end{solbox}


\subsection{4. 考点延展与防坑警示}
\begin{solbox}[分步推导与答题规范]
\textbf{[1]}  \textbf{极值视角（备选证法）}：
因 $f(x)$ 在闭区间 $[0, 1]$ 上连续，必取得最大值。由于 $f(0)=f(1)=0$ 且存在 $f(c)>0$，最大值必定在区间内部某点 $x_0 \in (0, 1)$ 处取得。\\ \\
由费马引理知 $f'(x_0) = 0$。若二阶导数 $f''(x) \ge 0$ 恒成立，则由凹凸性或泰勒公式 $f(x) \ge f(x_0) + f'(x_0)(x-x_0) = f(x_0)$，导出 $f(0) \ge f(x_0) > 0$，与 $f(0)=0$ 矛盾！故必存在点 $\xi$ 使得 $f''(\xi) < 0$。
\end{solbox}

---


\section{例题 11.3：分部积分法与夹逼定理求解含幂权积分极限}


\subsection{1. 题目描述}
计算极限：

\begin{equation*}
\lim\_{n \to \infty} \int\_0^1 (n+1) x^n \ln(1+x)\,\mathrm{d}x
\end{equation*}



\subsection{2. 解题思路与方法提炼}
\begin{solbox}[分步推导与答题规范]
\textbf{[1]}  \textbf{凑微分转化主要项}：
被积表达式中包含因子 $(n+1)x^n$，天然适合写成微分形式 $\mathrm{d}(x^{n+1})$。
\textbf{[2]}  \textbf{分部积分与余项控制}：
分部积分后，非积分项在端点直接求值；剩余的积分项由于包含 $\frac{x^{n+1}}{1+x}$，利用被积函数的有界性通过夹逼定理严格证明其余项极限为 $0$。
\end{solbox}


\subsection{3. 详细解答与规范步骤}
\begin{solbox}[分步推导与答题规范]
\textbf{[1]}  \textbf{执行分部积分}：

\begin{align*}
\begin{aligned}
I\_n \&= \int\_0^1 \ln(1+x)\,\mathrm{d}(x^{n+1}) \\ \\
\&= \left[ x^{n+1}\ln(1+x) \right]\_0^1 - \int\_0^1 x^{n+1} \cdot \frac{1}{1+x}\,\mathrm{d}x \\ \\
\&= 1^{n+1}\ln 2 - 0 - \int\_0^1 \frac{x^{n+1}}{1+x}\,\mathrm{d}x \\ \\
\&= \ln 2 - \int\_0^1 \frac{x^{n+1}}{1+x}\,\mathrm{d}x
\end{aligned}
\end{align*}


\textbf{[2]}  \textbf{对剩余积分项进行不等式夹逼}：
对任意 $x \in [0, 1]$，分母满足 $1 \le 1+x \le 2$，从而：

\begin{equation*}
0 \le \frac{x^{n+1}}{1+x} \le x^{n+1}
\end{equation*}

在区间 $[0, 1]$ 上取定积分：

\begin{equation*}
0 \le \int\_0^1 \frac{x^{n+1}}{1+x}\,\mathrm{d}x \le \int\_0^1 x^{n+1}\,\mathrm{d}x = \frac{1}{n+2}
\end{equation*}


\textbf{[3]}  \textbf{取极限得出最终结果}：
因为 $\lim_{n \to \infty} \frac{1}{n+2} = 0$，由夹逼定理得：

\begin{equation*}
\lim\_{n \to \infty} \int\_0^1 \frac{x^{n+1}}{1+x}\,\mathrm{d}x = 0
\end{equation*}

因此：

\begin{equation*}
\lim\_{n \to \infty} \int\_0^1 (n+1) x^n \ln(1+x)\,\mathrm{d}x = \ln 2 - 0 = \ln 2
\end{equation*}

\end{solbox}


\subsection{4. 考点延展与防坑警示}
\begin{solbox}[分步推导与答题规范]
\textbf{[1]}  \textbf{通式推广}：
若 $g(x)$ 在 $[0, 1]$ 上一阶连续可导，则有通用公式：

\begin{equation*}
\lim\_{n \to \infty} \int\_0^1 (n+1)x^n g(x)\,\mathrm{d}x = g(1)
\end{equation*}

该公式本质上体现了权函数 $(n+1)x^n$ 在 $n \to \infty$ 时集中在 $x=1$ 邻域的狄拉克 $\delta$ 脉冲特性。
\end{solbox}

---


\section{例题 11.4：黎曼引理变式：连续函数加权积分的渐近性质}


\subsection{1. 题目描述}
设函数 $f(x)$ 在闭区间 $[0, 1]$ 上连续。
\textbf{(1)}  证明：$\lim_{n \to \infty} \int_0^1 x^n f(x)\,\mathrm{d}x = 0$；
\textbf{(2)}  若进一步已知 $f(1) = 0$ 且 $f'(1)$ 存在，证明：$\lim_{n \to \infty} n \int_0^1 x^n f(x)\,\mathrm{d}x = 0$。


\subsection{2. 解题思路与方法提炼}
\begin{solbox}[分步推导与答题规范]
\textbf{[1]}  \textbf{有界性与积分直接放缩}：
由连续函数的有界性 $|f(x)| \le M$，可以直接控制被积函数，从而完成第 (1) 问证明。
\textbf{[2]}  \textbf{可导条件下的渐近阶分析}：
对于第 (2) 问，将 $n x^n$ 改写为 $\frac{n}{n+1}\mathrm{d}(x^{n+1})$ 后进行分部积分，结合条件 $f(1)=0$ 与导数定义完成极限估值。
\end{solbox}


\subsection{3. 详细解答与规范步骤}


\subsubsection{(1) 证明：$\lim_{n \to \infty} \int_0^1 x^n f(x)\,\mathrm{d}x = 0$}
\begin{solbox}[分步推导与答题规范]
\textbf{[1]}  \textbf{利用极值定理确界}：
由于 $f(x)$ 在 $[0, 1]$ 上连续，存在常数 $M > 0$，使得对一切 $x \in [0, 1]$ 均有 $|f(x)| \le M$。

\textbf{[2]}  \textbf{绝对值不等式放缩}：

\begin{equation*}
\left| \int\_0^1 x^n f(x)\,\mathrm{d}x \right| \le \int\_0^1 x^n |f(x)|\,\mathrm{d}x \le M \int\_0^1 x^n\,\mathrm{d}x = \frac{M}{n+1}
\end{equation*}

因为 $\lim_{n \to \infty} \frac{M}{n+1} = 0$，由夹逼定理立即证得：

\begin{equation*}
\lim\_{n \to \infty} \int\_0^1 x^n f(x)\,\mathrm{d}x = 0
\end{equation*}

\end{solbox}


\subsubsection{(2) 证明：$\lim_{n \to \infty} n \int_0^1 x^n f(x)\,\mathrm{d}x = 0$}
\begin{solbox}[分步推导与答题规范]
\textbf{[1]}  \textbf{分部积分转化形式}：

\begin{align*}
\begin{aligned}
\int\_0^1 x^n f(x)\,\mathrm{d}x \&= \frac{1}{n+1}\int\_0^1 f(x)\,\mathrm{d}(x^{n+1}) \\ \\
\&= \frac{1}{n+1}\left[ x^{n+1}f(x) \right]\_0^1 - \frac{1}{n+1}\int\_0^1 x^{n+1}f'(x)\,\mathrm{d}x \quad (\text{假设 } f' \text{ 连续})
\end{aligned}
\end{align*}

若 $f'(x)$ 不连续，只需利用导数定义：因 $f(1)=0$ 且 $f'(1)$ 存在，对任意 $\varepsilon > 0$，存在 $\delta \in (0, 1)$，当 $x \in [1-\delta, 1]$ 时：

\begin{equation*}
|f(x)| = |f(x) - f(1)| \le (|f'(1)| + 1)(1-x)
\end{equation*}


\textbf{[2]}  \textbf{区间分段估计}：
将积分分为 $[0, 1-\delta]$ 与 $[1-\delta, 1]$ 两部分：
在 $[0, 1-\delta]$ 上：

\begin{equation*}
n \left| \int\_0^{1-\delta} x^n f(x)\,\mathrm{d}x \right| \le n M \int\_0^{1-\delta} x^n\,\mathrm{d}x = \frac{n M (1-\delta)^{n+1}}{n+1} \to 0 \quad (n \to \infty)
\end{equation*}

在 $[1-\delta, 1]$ 上：

\begin{align*}
\begin{aligned}
n \left| \int\_{1-\delta}^1 x^n f(x)\,\mathrm{d}x \right| \&\le n (|f'(1)|+1) \int\_{1-\delta}^1 x^n (1-x)\,\mathrm{d}x \\ \\
\&= n (|f'(1)|+1) \left( \frac{1-(1-\delta)^{n+1}}{n+1} - \frac{1-(1-\delta)^{n+2}}{n+2} \right) \\ \\
\&< n (|f'(1)|+1) \left( \frac{1}{n+1} - \frac{1}{n+2} \right) = \frac{n (|f'(1)|+1)}{(n+1)(n+2)} \to 0
\end{aligned}
\end{align*}

综上，取极限即证得：

\begin{equation*}
\lim\_{n \to \infty} n \int\_0^1 x^n f(x)\,\mathrm{d}x = 0
\end{equation*}

\end{solbox}


\subsection{4. 考点延展与防坑警示}
\begin{solbox}[分步推导与答题规范]
\textbf{[1]}  注意此处的阶数关系：若 $f(1) \ne 0$，则 $\lim_{n \to \infty} n \int_0^1 x^n f(x)\,\mathrm{d}x = f(1) \ne 0$。只有当 $f(1)=0$ 时，高一阶的极限才为 $0$。
\end{solbox}

---


\section{例题 11.5：锯齿波函数的均值极限与周期性定积分估计}


\subsection{1. 题目描述}
设锯齿波函数 $f(x) = x - [x]$，其中 $[x]$ 表示不大于 $x$ 的最大整数。
计算极限：

\begin{equation*}
\lim\_{x \to +\infty} \frac{1}{x} \int\_0^x f(t)\,\mathrm{d}t
\end{equation*}



\subsection{2. 解题思路与方法提炼}
\begin{solbox}[分步推导与答题规范]
\textbf{[1]}  \textbf{识别周期函数及其单周期积分}：
函数 $f(t) = t - [t]$ 是以 $T=1$ 为基本周期的周期函数。计算单个周期内的积分值 $\int_0^1 t\,\mathrm{d}t = \frac{1}{2}$。
\textbf{[2]}  \textbf{利用整数分段与夹逼定理求解无界极限}：
令 $n = [x]$，则 $n \le x < n+1$。将积分区间 $[0, x]$ 分解为 $n$ 个完整周期 $[0, n]$ 与不足一个周期的剩余部分 $[n, x]$，通过不等式双向夹逼求出平均值。
\end{solbox}


\subsection{3. 详细解答与规范步骤}
\begin{solbox}[分步推导与答题规范]
\textbf{[1]}  \textbf{单周期积分计算}：
对于任意整数 $k$，当 $t \in [k, k+1)$ 时，$[t] = k$，故 $f(t) = t - k$。
在每个单位区间上的定积分为：

\begin{equation*}
\int\_k^{k+1} f(t)\,\mathrm{d}t = \int\_k^{k+1} (t-k)\,\mathrm{d}t = \int\_0^1 u\,\mathrm{d}u = \frac{1}{2}
\end{equation*}


\textbf{[2]}  \textit{}对任意正实数 $x$ 进行分段分解\textit{}：
设 $n = [x] \in \mathbb{N}$，则有：

\begin{equation*}
n \le x < n+1
\end{equation*}

将定积分分解为：

\begin{equation*}
\int\_0^x f(t)\,\mathrm{d}t = \int\_0^n f(t)\,\mathrm{d}t + \int\_n^x f(t)\,\mathrm{d}t
\end{equation*}

其中完整周期的积分和为：

\begin{equation*}
\int\_0^n f(t)\,\mathrm{d}t = \sum\_{k=0}^{n-1} \int\_k^{k+1} f(t)\,\mathrm{d}t = \sum\_{k=0}^{n-1} \frac{1}{2} = \frac{n}{2}
\end{equation*}

对于剩余部分，由于对一切 $t$ 恒有 $0 \le f(t) < 1$，故：

\begin{equation*}
0 \le \int\_n^x f(t)\,\mathrm{d}t < \int\_n^{n+1} 1\,\mathrm{d}t = 1
\end{equation*}


\textbf{[3]}  \textbf{构建不等式夹逼}：
综合上述结果，原积分满足：

\begin{equation*}
\frac{n}{2} \le \int\_0^x f(t)\,\mathrm{d}t < \frac{n}{2} + 1
\end{equation*}

两端同时除以 $x$（注意 $n \le x < n+1$）：

\begin{equation*}
\frac{n}{2(n+1)} < \frac{n}{2x} \le \frac{1}{x}\int\_0^x f(t)\,\mathrm{d}t < \frac{\frac{n}{2} + 1}{n} = \frac{1}{2} + \frac{1}{n}
\end{equation*}

当 $x \to +\infty$ 时，必有 $n = [x] \to +\infty$。\\ \\
左侧极限为 $\lim_{n\to\infty} \frac{n}{2(n+1)} = \frac{1}{2}$，右侧极限为 $\lim_{n\to\infty} \left(\frac{1}{2} + \frac{1}{n}\right) = \frac{1}{2}$。\\ \\
由夹逼定理得：

\begin{equation*}
\lim\_{x \to +\infty} \frac{1}{x} \int\_0^x f(t)\,\mathrm{d}t = \frac{1}{2}
\end{equation*}

\end{solbox}


\subsection{4. 考点延展与防坑警示}
\begin{solbox}[分步推导与答题规范]
\textbf{[1]}  \textbf{一般周期函数的无穷均值定理}：
若 $g(t)$ 是以 $T$ 为周期的连续函数，则恒有：

\begin{equation*}
\lim\_{x \to +\infty} \frac{1}{x}\int\_0^x g(t)\,\mathrm{d}t = \frac{1}{T}\int\_0^T g(t)\,\mathrm{d}t
\end{equation*}

本题正是该定理在非连续分段周期函数上的经典推导范式。
\end{solbox}

---


\section{例题 11.6：对称二次分部积分与二阶导数极值放缩}


\subsection{1. 题目描述}
设函数 $f(x)$ 在闭区间 $[0, 1]$ 上具有连续的二阶导数，且满足端点值条件：

\begin{equation*}
f(0) = f(1) = 0
\end{equation*}

设 $M = \max_{x \in [0, 1]} |f''(x)|$。证明：

\begin{equation*}
\left| \int\_0^1 f(x)\,\mathrm{d}x \right| \le \frac{1}{12} M
\end{equation*}



\subsection{2. 解题思路与方法提炼}
\begin{solbox}[分步推导与答题规范]
\textbf{[1]}  \textbf{构造二次对称核函数}：
直接分部积分 $\int_0^1 f(x)\,\mathrm{d}x$ 会产生一阶导数项，但题目给出的限制是二阶导数 $M$。因此需要进行两次分部积分，将积分核提升为二次多项式 $P(x)$。
\textbf{[2]}  \textbf{消除边界项的核函数配凑}：
为使边界项 $[P(x)f'(x)]_0^1$ 与 $[P'(x)f(x)]_0^1$ 全部消去，取 $P(x) = \frac{1}{2}x(x-1)$。此时 $P''(x) = 1$，$P(0)=P(1)=0$，从而实现端点项全消。
\end{solbox}


\subsection{3. 详细解答与规范步骤}
\begin{solbox}[分步推导与答题规范]
\textbf{[1]}  \textbf{验证积分恒等式}：
设核函数 $P(x) = \frac{1}{2}x(x-1) = \frac{1}{2}(x^2 - x)$，则其一阶导数与二阶导数分别为：

\begin{equation*}
P'(x) = x - \frac{1}{2},\quad P''(x) = 1
\end{equation*}

考虑积分 $\int_0^1 P(x)f''(x)\,\mathrm{d}x$，进行分部积分：

\begin{align*}
\begin{aligned}
\int\_0^1 P(x)f''(x)\,\mathrm{d}x \&= \left[ P(x)f'(x) \right]\_0^1 - \int\_0^1 P'(x)f'(x)\,\mathrm{d}x
\end{aligned}
\end{align*}

由于 $P(0) = P(1) = 0$，首项 $\left[ P(x)f'(x) \right]_0^1 = 0$。
继续对第二项进行分部积分：

\begin{align*}
\begin{aligned}
-\int\_0^1 P'(x)f'(x)\,\mathrm{d}x \&= -\left[ P'(x)f(x) \right]\_0^1 + \int\_0^1 P''(x)f(x)\,\mathrm{d}x
\end{aligned}
\end{align*}

已知 $f(0) = f(1) = 0$，故边界项 $-\left[ P'(x)f(x) \right]_0^1 = 0$。\\ \\
又因 $P''(x) = 1$，上式化简为：

\begin{equation*}
\int\_0^1 P(x)f''(x)\,\mathrm{d}x = \int\_0^1 f(x)\,\mathrm{d}x
\end{equation*}

即恒有精确积分等式：

\begin{equation*}
\int\_0^1 f(x)\,\mathrm{d}x = \frac{1}{2}\int\_0^1 x(x-1)f''(x)\,\mathrm{d}x
\end{equation*}


\textbf{[2]}  \textbf{绝对值不等式精确放缩}：
两端取绝对值并应用定积分不等式：

\begin{align*}
\begin{aligned}
\left| \int\_0^1 f(x)\,\mathrm{d}x \right| \&= \left| \frac{1}{2}\int\_0^1 x(x-1)f''(x)\,\mathrm{d}x \right| \\ \\
\&\le \frac{1}{2}\int\_0^1 |x(x-1)| \cdot |f''(x)|\,\mathrm{d}x \\ \\
\&\le \frac{M}{2} \int\_0^1 x(1-x)\,\mathrm{d}x
\end{aligned}
\end{align*}


\textbf{[3]}  \textbf{定积分定值计算}：
计算定积分：

\begin{equation*}
\int\_0^1 x(1-x)\,\mathrm{d}x = \int\_0^1 (x - x^2)\,\mathrm{d}x = \left[ \frac{x^2}{2} - \frac{x^3}{3} \right]\_0^1 = \frac{1}{2} - \frac{1}{3} = \frac{1}{6}
\end{equation*}

代入得：

\begin{equation*}
\left| \int\_0^1 f(x)\,\mathrm{d}x \right| \le \frac{M}{2} \cdot \frac{1}{6} = \frac{1}{12}M
\end{equation*}

命题得证。
\end{solbox}


\subsection{4. 考点延展与防坑警示}
\begin{solbox}[分步推导与答题规范]
\textbf{[1]}  \textbf{系数的最优性}：常数 $\frac{1}{12}$ 是无法改进的最优常数。当 $f(x) = \frac{M}{2}x(1-x)$ 时，$f''(x) = -M$，此时 $|\int_0^1 f(x)\mathrm{d}x| = \frac{M}{2} \cdot \frac{1}{6} = \frac{1}{12}M$，不等式等号恰好成立。
\end{solbox}

---


\section{例题 11.7：变上限辅助函数单调性法证明积分不等式}


\subsection{1. 题目描述}
设函数 $f(x)$ 在闭区间 $[0, 1]$ 上连续且单调不减。证明：

\begin{equation*}
\int\_0^1 x f(x)\,\mathrm{d}x \ge \frac{1}{2}\int\_0^1 f(x)\,\mathrm{d}x
\end{equation*}



\subsection{2. 解题思路与方法提炼}
\begin{solbox}[分步推导与答题规范]
\textbf{[1]}  \textbf{构造变上限积分辅助函数}：
将常数上限 $1$ 替换为自变量 $t \in [0, 1]$，构造辅助函数：

\begin{equation*}
F(t) = \int\_0^t x f(x)\,\mathrm{d}x - \frac{t}{2}\int\_0^t f(x)\,\mathrm{d}x
\end{equation*}

证明 $F'(t) \ge 0$，结合 $F(0)=0$ 即可推出 $F(1) \ge 0$。
\textbf{[2]}  \textbf{对称化区间重积分视角（二重积分法）}：
将一维积分乘积化为二重对称积分，利用 $(x-y)(f(x)-f(y)) \ge 0$ 的单调性代数结构瞬时获证。
\end{solbox}


\subsection{3. 详细解答与规范步骤}


\subsubsection{方法一：变上限积分辅助函数求导法}
\begin{solbox}[分步推导与答题规范]
\textbf{[1]}  \textbf{设定辅助函数}：
构造辅助函数 $F(t)$，定义于 $t \in [0, 1]$：

\begin{equation*}
F(t) = \int\_0^t x f(x)\,\mathrm{d}x - \frac{t}{2}\int\_0^t f(x)\,\mathrm{d}x
\end{equation*}

显然 $F(0) = 0$。

\textbf{[2]}  \textbf{对辅助函数求导}：
由变上限积分求导法则：

\begin{align*}
\begin{aligned}
F'(t) \&= t f(t) - \left( \frac{1}{2}\int\_0^t f(x)\,\mathrm{d}x + \frac{t}{2}f(t) \right) \\ \\
\&= \frac{t}{2}f(t) - \frac{1}{2}\int\_0^t f(x)\,\mathrm{d}x \\ \\
\&= \frac{1}{2}\int\_0^t [f(t) - f(x)]\,\mathrm{d}x
\end{aligned}
\end{align*}


\textbf{[3]}  \textbf{利用单调性判断导数正负}：
因为 $f(x)$ 在 $[0, 1]$ 上单调不减，所以当 $x \in [0, t]$ 时，恒有 $f(t) \ge f(x)$，即：

\begin{equation*}
f(t) - f(x) \ge 0
\end{equation*}

因此被积函数非负，定积分 $F'(t) \ge 0$。\\ \\
这表明 $F(t)$ 在 $[0, 1]$ 上单调递增。由 $F(0) = 0$ 可得：

\begin{equation*}
F(1) \ge F(0) = 0
\end{equation*}

即：

\begin{equation*}
\int\_0^1 x f(x)\,\mathrm{d}x - \frac{1}{2}\int\_0^1 f(x)\,\mathrm{d}x \ge 0 \implies \int\_0^1 x f(x)\,\mathrm{d}x \ge \frac{1}{2}\int\_0^1 f(x)\,\mathrm{d}x
\end{equation*}

\end{solbox}


\subsubsection{方法二：二重积分对称化法（最美秒杀视角）}
\begin{solbox}[分步推导与答题规范]
\textbf{[1]}  \textbf{构造对称二重积分}：
考虑矩形区域 $D = [0, 1] \times [0, 1]$ 上的二重积分：

\begin{equation*}
I = \iint\_D (x - y)[f(x) - f(y)]\,\mathrm{d}x\mathrm{d}y
\end{equation*}

由于 $f$ 单调不减，$x - y$ 与 $f(x) - f(y)$ 符号始终同向，故 $(x - y)[f(x) - f(y)] \ge 0$ 恒成立。因此：

\begin{equation*}
I \ge 0
\end{equation*}


\textbf{[2]}  \textbf{展开并利用变量对称性}：

\begin{align*}
\begin{aligned}
I \&= \iint\_D [x f(x) - x f(y) - y f(x) + y f(y)]\,\mathrm{d}x\mathrm{d}y \\ \\
\&= \int\_0^1 x f(x)\,\mathrm{d}x \int\_0^1 1\,\mathrm{d}y - \int\_0^1 x\,\mathrm{d}x \int\_0^1 f(y)\,\mathrm{d}y - \int\_0^1 y\,\mathrm{d}y \int\_0^1 f(x)\,\mathrm{d}x + \int\_0^1 1\,\mathrm{d}x \int\_0^1 y f(y)\,\mathrm{d}y \\ \\
\&= 2 \int\_0^1 x f(x)\,\mathrm{d}x - 2 \left( \frac{1}{2} \right) \int\_0^1 f(x)\,\mathrm{d}x \\ \\
\&= 2 \left( \int\_0^1 x f(x)\,\mathrm{d}x - \frac{1}{2}\int\_0^1 f(x)\,\mathrm{d}x \right)
\end{aligned}
\end{align*}

因为 $I \ge 0$，立即得出：

\begin{equation*}
\int\_0^1 x f(x)\,\mathrm{d}x \ge \frac{1}{2}\int\_0^1 f(x)\,\mathrm{d}x
\end{equation*}

\end{solbox}


\subsection{4. 考点延展与防坑警示}
\begin{solbox}[分步推导与答题规范]
\textbf{[1]}  二重积分对称化方法是证明切比雪夫积分不等式（Chebyshev's integral inequality）的标准数学引擎。考研大题中若使用方法二，建议写出累次积分展开过程，逻辑严丝合缝且具备极高的学术审美。
\end{solbox}

---


\section{例题 11.8：一阶导数有界条件下中值分割的积分估计}


\subsection{1. 题目描述}
设函数 $f(x)$ 在闭区间 $[0, 1]$ 上可导，满足端点条件 $f(0) = f(1) = 0$，且导函数满足 $|f'(x)| \le M$。
证明：

\begin{equation*}
\left| \int\_0^1 f(x)\,\mathrm{d}x \right| \le \frac{1}{4} M
\end{equation*}



\subsection{2. 解题思路与方法提炼}
\begin{solbox}[分步推导与答题规范]
\textbf{[1]}  \textbf{区间双向拉格朗日中值放缩}：
对区间内任意一点 $x \in [0, 1]$，既可以由左端点 $0$ 展开 $f(x) - f(0) = f'(\xi_1)x$，也可以由右端点 $1$ 展开 $f(x) - f(1) = f'(\xi_2)(x-1)$。
\textbf{[2]}  \textbf{折线外包络线控制}：
由导数有界性可知 $|f(x)| \le M x$ 且 $|f(x)| \le M(1-x)$，故 $|f(x)| \le M \min\{x, 1-x\}$。在对称中心 $x = \frac{1}{2}$ 处分段积分即可求得 $\frac{1}{4}M$。
\end{solbox}


\subsection{3. 详细解答与规范步骤}
\begin{solbox}[分步推导与答题规范]
\textbf{[1]}  \textbf{双端点微分中值放缩}：
对任意给定的 $x \in (0, 1)$：\\ \\
在区间 $[0, x]$ 上应用拉格朗日中值定理：

\begin{equation*}
f(x) = f(x) - f(0) = f'(\xi\_1) x \quad (\xi\_1 \in (0, x))
\end{equation*}

因此：

\begin{equation*}
|f(x)| \le |f'(\xi\_1)| x \le M x
\end{equation*}

同理，在区间 $[x, 1]$ 上应用拉格朗日中值定理：

\begin{equation*}
f(x) = f(x) - f(1) = f'(\xi\_2) (x - 1) \quad (\xi\_2 \in (x, 1))
\end{equation*}

因此：

\begin{equation*}
|f(x)| \le |f'(\xi\_2)| (1 - x) \le M (1 - x)
\end{equation*}


\textbf{[2]}  \textbf{取包络最小值构建分段函数}：
由此可得，对一切 $x \in [0, 1]$，均有：

\begin{equation*}
|f(x)| \le M \min\{x, 1-x\}
\end{equation*}


\textbf{[3]}  \textbf{分段积分计算上界}：
应用定积分绝对值不等式：

\begin{align*}
\begin{aligned}
\left| \int\_0^1 f(x)\,\mathrm{d}x \right| \&\le \int\_0^1 |f(x)|\,\mathrm{d}x \\ \\
\&\le M \int\_0^1 \min\{x, 1-x\}\,\mathrm{d}x \\ \\
\&= M \left( \int\_0^{1/2} x\,\mathrm{d}x + \int\_{1/2}^1 (1-x)\,\mathrm{d}x \right) \\ \\
\&= M \left( \left[ \frac{x^2}{2} \right]\_0^{1/2} + \left[ x - \frac{x^2}{2} \right]\_{1/2}^1 \right) \\ \\
\&= M \left( \frac{1}{8} + \frac{1}{8} \right) = \frac{1}{4}M
\end{aligned}
\end{align*}

命题得证。
\end{solbox}


\subsection{4. 考点延展与防坑警示}
\begin{solbox}[分步推导与答题规范]
\textbf{[1]}  \textbf{若只有一个端点为零}：若仅已知 $f(0)=0$，则 $|f(x)| \le Mx$，积分上界为 $\int_0^1 Mx\,\mathrm{d}x = \frac{1}{2}M$。
\textbf{[2]}  \textbf{两端点均为零时}，务必利用几何对称性在 $x = 1/2$ 处折半取交集，获得更为紧凑紧致的 $\frac{1}{4}M$ 上界。
\end{solbox}

---


\section{例题 11.9：二阶导数有界与内点泰勒展开积分估计}


\subsection{1. 题目描述}
设函数 $f(x)$ 在闭区间 $[0, 2]$ 上二阶可导，且满足在区间中点处 $f(1) = 0, f'(1) = 0$。
若对一切 $x \in [0, 2]$ 均有 $|f''(x)| \le M$。证明：

\begin{equation*}
\left| \int\_0^2 f(x)\,\mathrm{d}x \right| \le \frac{1}{3} M
\end{equation*}



\subsection{2. 解题思路与方法提炼}
\begin{solbox}[分步推导与答题规范]
\textbf{[1]}  \textbf{选择泰勒展开中心}：
由于已知中点 $x_0 = 1$ 处函数值与一阶导数值均为 $0$，以 $x_0 = 1$ 为中心展开能最大限度消除低阶多项式项。
\textbf{[2]}  \textbf{拉格朗日余项放缩积分}：
直接写出带拉格朗日余项的一阶泰勒公式，被积函数被二次多项式 $\frac{M}{2}(x-1)^2$ 紧致控制，积分立得结果。
\end{solbox}


\subsection{3. 详细解答与规范步骤}
\begin{solbox}[分步推导与答题规范]
\textbf{[1]}  \textit{}在 $x_0 = 1$ 处写出泰勒展开式\textit{}：
由泰勒定理，对任意 $x \in [0, 2]$，将 $f(x)$ 在 $x_0 = 1$ 处展开至一阶（带二阶拉格朗日余项）：

\begin{equation*}
f(x) = f(1) + f'(1)(x - 1) + \frac{f''(\xi)}{2}(x - 1)^2 \quad (\xi \text{ 介于 } 1 \text{ 与 } x \text{ 之间})
\end{equation*}

已知 $f(1) = 0$ 且 $f'(1) = 0$，上式化简为：

\begin{equation*}
f(x) = \frac{f''(\xi)}{2}(x - 1)^2
\end{equation*}


\textbf{[2]}  \textbf{取绝对值不等式逐项放缩}：
由题设 $|f''(x)| \le M$，故对任意 $\xi \in [0, 2]$ 恒有 $|f''(\xi)| \le M$。因此：

\begin{equation*}
|f(x)| = \left| \frac{f''(\xi)}{2}(x - 1)^2 \right| \le \frac{M}{2}(x - 1)^2
\end{equation*}


\textbf{[3]}  \textbf{逐项定积分计算}：
利用定积分绝对值不等式：

\begin{align*}
\begin{aligned}
\left| \int\_0^2 f(x)\,\mathrm{d}x \right| \&\le \int\_0^2 |f(x)|\,\mathrm{d}x \\ \\
\&\le \frac{M}{2} \int\_0^2 (x - 1)^2\,\mathrm{d}x \\ \\
\&= \frac{M}{2} \left[ \frac{(x - 1)^3}{3} \right]\_0^2 \\ \\
\&= \frac{M}{2} \left( \frac{1^3}{3} - \frac{(-1)^3}{3} \right) = \frac{M}{2} \cdot \frac{2}{3} = \frac{1}{3}M
\end{aligned}
\end{align*}

命题得证。
\end{solbox}


\subsection{4. 考点延展与防坑警示}
\begin{solbox}[分步推导与答题规范]
\textbf{[1]}  \textbf{展开中心选择原则}：泰勒公式在积分不等式证明中，展开中心务必选择\textbf{已知导数信息最多、对称性最好的点}。若在本题选择端点 $0$ 展开，由于 $f(0), f'(0)$ 均未知，将陷入无法消去低阶常数的死胡同。
\end{solbox}

---


\section{例题 11.10：单调函数高频振荡积分的衰减阶估计}


\subsection{1. 题目描述}
设函数 $f(x)$ 在闭区间 $[0, 2\pi]$ 上单调递增且具有连续的一阶导数。证明：对任意正整数 $n$，均有：

\begin{equation*}
\left| \int\_0^{2\pi} f(x)\sin nx\,\mathrm{d}x \right| \le \frac{2}{n} [f(2\pi) - f(0)]
\end{equation*}



\subsection{2. 解题思路与方法提炼}
\begin{solbox}[分步推导与答题规范]
\textbf{[1]}  \textbf{分部积分抽取振荡因子}：
利用 $\sin nx\,\mathrm{d}x = -\frac{1}{n}\mathrm{d}(\cos nx)$ 进行分部积分，将振荡频率转化为幅度衰减因子 $\frac{1}{n}$。
\textbf{[2]}  \textbf{利用单调性消除绝对值符号}：
分部积分后的剩余积分中含有 $f'(x)$。由于 $f(x)$ 单调递增，导数 $f'(x) \ge 0$ 恒成立，因而绝对值可以与导数完全解耦。
\end{solbox}


\subsection{3. 详细解答与规范步骤}
\begin{solbox}[分步推导与答题规范]
\textbf{[1]}  \textbf{执行分部积分}：

\begin{align*}
\begin{aligned}
\int\_0^{2\pi} f(x)\sin nx\,\mathrm{d}x \&= -\frac{1}{n}\int\_0^{2\pi} f(x)\,\mathrm{d}(\cos nx) \\ \\
\&= -\frac{1}{n}\left[ f(x)\cos nx \right]\_0^{2\pi} + \frac{1}{n}\int\_0^{2\pi} f'(x)\cos nx\,\mathrm{d}x
\end{aligned}
\end{align*}


\textbf{[2]}  \textbf{计算端点项}：
因为 $\cos(2n\pi) = 1, \cos(0) = 1$，首项为：

\begin{equation*}
-\frac{1}{n}[f(2\pi) \cdot 1 - f(0) \cdot 1] = -\frac{1}{n}[f(2\pi) - f(0)]
\end{equation*}

从而：

\begin{equation*}
\int\_0^{2\pi} f(x)\sin nx\,\mathrm{d}x = -\frac{1}{n}[f(2\pi) - f(0)] + \frac{1}{n}\int\_0^{2\pi} f'(x)\cos nx\,\mathrm{d}x
\end{equation*}


\textbf{[3]}  \textbf{三角绝对值不等式放缩}：
两端取绝对值：

\begin{align*}
\begin{aligned}
\left| \int\_0^{2\pi} f(x)\sin nx\,\mathrm{d}x \right| \&\le \frac{1}{n}|f(2\pi) - f(0)| + \frac{1}{n}\left| \int\_0^{2\pi} f'(x)\cos nx\,\mathrm{d}x \right| \\ \\
\&\le \frac{1}{n}[f(2\pi) - f(0)] + \frac{1}{n}\int\_0^{2\pi} |f'(x)| \cdot |\cos nx|\,\mathrm{d}x
\end{aligned}
\end{align*}

由于 $|\cos nx| \le 1$，且已知 $f(x)$ 单调递增即 $f'(x) \ge 0 \implies |f'(x)| = f'(x)$：

\begin{equation*}
\int\_0^{2\pi} |f'(x)| \cdot |\cos nx|\,\mathrm{d}x \le \int\_0^{2\pi} f'(x)\,\mathrm{d}x = f(2\pi) - f(0)
\end{equation*}


\textbf{[4]}  \textbf{合并上界}：
将上述估计代入原式：

\begin{equation*}
\left| \int\_0^{2\pi} f(x)\sin nx\,\mathrm{d}x \right| \le \frac{1}{n}[f(2\pi) - f(0)] + \frac{1}{n}[f(2\pi) - f(0)] = \frac{2}{n}[f(2\pi) - f(0)]
\end{equation*}

命题得证。
\end{solbox}


\subsection{4. 考点延展与防坑警示}
\begin{solbox}[分步推导与答题规范]
\textbf{[1]}  \textbf{黎曼-勒贝格引理的定量形式}：该结论直接给出了狄利克雷核/傅里叶积分衰减速率 $O(1/n)$ 的显式常数估计。
\textbf{[2]}  \textbf{第二积分中值定理秒杀视角}：由积分第二中值定理，存在 $\xi \in [0, 2\pi]$ 使得 $\int_0^{2\pi} f(x)\sin nx\,\mathrm{d}x = f(0)\int_0^\xi \sin nx\,\mathrm{d}x + f(2\pi)\int_\xi^{2\pi} \sin nx\,\mathrm{d}x$。利用 $\left|\int_\alpha^\beta \sin nx\,\mathrm{d}x\right| \le \frac{2}{n}$ 亦可快速获证。
\end{solbox}

---


\section{例题 11.11：双向微积分基本定理与连续模全域控制}


\subsection{1. 题目描述}
设函数 $f(x)$ 在闭区间 $[a, b]$ 上具有一阶连续导数，且满足端点条件：

\begin{equation*}
f(a) = f(b) = 0
\end{equation*}

证明：对任意 $x \in [a, b]$，均有：

\begin{equation*}
|f(x)| \le \frac{1}{2}\int\_a^b |f'(t)|\,\mathrm{d}t
\end{equation*}



\subsection{2. 解题思路与方法提炼}
\begin{solbox}[分步推导与答题规范]
\textbf{[1]}  \textbf{双端向内应用微积分基本定理（牛顿-莱布尼茨公式）}：
对于区间内任意一点 $x$，既可以写成从左端点 $a$ 出发的积分，也可以写成从右端点 $b$ 出发的积分：

\begin{equation*}
f(x) = \int\_a^x f'(t)\,\mathrm{d}t,\quad f(x) = -\int\_x^b f'(t)\,\mathrm{d}t
\end{equation*}

\textbf{[2]}  \textbf{算术平均与绝对值不等式}：
将两式取绝对值后相加并除以 $2$，由于 $[a, x]$ 与 $[x, b]$ 恰好无缝拼合为全区间 $[a, b]$，直接得到系数 $\frac{1}{2}$。
\end{solbox}


\subsection{3. 详细解答与规范步骤}
\begin{solbox}[分步推导与答题规范]
\textbf{[1]}  \textbf{双向表示函数值}：
由牛顿-莱布尼茨公式及条件 $f(a) = 0$：

\begin{equation*}
f(x) = f(x) - f(a) = \int\_a^x f'(t)\,\mathrm{d}t
\end{equation*}

同理，结合条件 $f(b) = 0$：

\begin{equation*}
f(x) = f(x) - f(b) = -\int\_x^b f'(t)\,\mathrm{d}t
\end{equation*}


\textbf{[2]}  \textbf{取绝对值放缩}：
对第一式两边取绝对值：

\begin{equation*}
|f(x)| = \left| \int\_a^x f'(t)\,\mathrm{d}t \right| \le \int\_a^x |f'(t)|\,\mathrm{d}t
\end{equation*}

对第二式两边取绝对值：

\begin{equation*}
|f(x)| = \left| -\int\_x^b f'(t)\,\mathrm{d}t \right| \le \int\_x^b |f'(t)|\,\mathrm{d}t
\end{equation*}


\textbf{[3]}  \textbf{相加取平均完成证明}：
将两式不等式直接相加：

\begin{equation*}
2|f(x)| \le \int\_a^x |f'(t)|\,\mathrm{d}t + \int\_x^b |f'(t)|\,\mathrm{d}t
\end{equation*}

由定积分对区间的可加性，右端两项合并为：

\begin{equation*}
\int\_a^x |f'(t)|\,\mathrm{d}t + \int\_x^b |f'(t)|\,\mathrm{d}t = \int\_a^b |f'(t)|\,\mathrm{d}t
\end{equation*}

两边同除以 $2$，即得：

\begin{equation*}
|f(x)| \le \frac{1}{2}\int\_a^b |f'(t)|\,\mathrm{d}t
\end{equation*}

由于 $x \in [a, b]$ 为任意取定，不等式在全区间一致成立。
\end{solbox}


\subsection{4. 考点延展与防坑警示}
\begin{solbox}[分步推导与答题规范]
\textbf{[1]}  \textbf{庞加莱（Poincaré）型不等式的雏形}：两端平方并结合柯西-施瓦茨不等式，可进一步导出 $L^2$ 范数下的经典庞加莱不等式 $\int_a^b |f(x)|^2\,\mathrm{d}x \le C (b-a)^2 \int_a^b |f'(x)|^2\,\mathrm{d}x$。
\textbf{[2]}  \textbf{系数的最优性检验}：若函数 $f(x)$ 在中点两边具有斜率相等的对称折线形状（或在光滑逼近下），峰值点 $|f(x_0)|$ 恰好等于总变差的一半，故常数 $\frac{1}{2}$ 是严格不可优化的最优界。
\end{solbox}

\end{document}
